盈利改善是真实的，但“利润高增”与“利润质量改善”不是同一个结论。2025年同一控制下重述口径显示收入增长13.97\%、归母增长86.00\%，同时经营现金流由重述2024年的147.27亿元降至77.67亿元。2026Q1利润与现金同步较强，H1却只有利润预告；正式半年报必须补齐收入、毛利、现金与营运资金。

\section{E04｜2024重述至2026H1：利润弹性强，口径必须一致}

\begin{exhibitbox}[合并口径财务进展]
\small
\begin{tabularx}{\exhibitboxwidth}{L{2.8cm}L{3.0cm}R{2.2cm}R{2.2cm}R{2.2cm}X}
\toprule
\textbf{期间} & \textbf{口径} & \textbf{收入} & \textbf{归母} & \textbf{扣非} & \textbf{OCF} \\
\midrule
2024A & 同控重述比较期 & 1,333.51 & 42.20 & 30.72 & 147.27 \\
2025A & 持续经营实际 & 1,519.78 & 78.48 & 61.26 & 77.67 \\
2026Q1A & 未经审计季报 & 433.12 & 48.32 & 47.67 & 81.52 \\
2026H1预告 & 未经审计区间 & \na & 92--110 & 90--108 & \na \\
\bottomrule
\end{tabularx}
\sourcenote{CF-03、CF-04、CF-05。金额单位：亿元；H1为业绩预告性质。}
\end{exhibitbox}

2025归母增长快于收入，说明利润率和经营结构改善；扣非利润也接近翻倍，不能简单归因于一次性收益。但OCF/归母由2024重述的约3.49倍降至2025的约0.99倍，反映生产爬坡与营运资金投入。2026Q1 OCF 81.52亿元高于归母48.32亿元，单季改善仍需全年节奏验证。

\section{核心业务毛利是高价订单兑现的最短路径}

2025年船舶造修及海工收入1,312.79亿元、毛利率11.72\%，主营业务毛利率12.27\%。相对2024，核心毛利率提升2.94个百分点，说明价格、结构和效率开始兑现。未来最重要的并非单季合并毛利率，而是可比口径核心业务毛利能否在设备涨价和产量爬坡下继续改善。\srcid{CF-03}

\begin{exhibitbox}[2025收入与利润锚]
\small
\begin{tabularx}{\exhibitboxwidth}{L{4.2cm}R{2.8cm}L{3.5cm}X}
\toprule
\textbf{指标} & \textbf{数值} & \textbf{模型用途} & \textbf{边界} \\
\midrule
合并营业收入 & 1,519.78亿元 & 总收入基期 & 同控重述持续经营口径 \\
造船/修船/海工收入 & 1,312.79亿元 & Growth业务基期 & 含多种业务和确认方式 \\
配套、机电及其他主营收入 & 186.17亿元 & Base业务主要组成 & 独立毛利未完整披露 \\
合并毛利率 & 12.58\% & 合并勾稽 & 不等同核心业务毛利 \\
核心业务毛利率 & 11.72\% & 高价订单兑现锚 & 不外推为船型毛利 \\
归母净利润 & 78.48亿元 & 盈利基期 & 2025 EPS使用加权分母 \\
\bottomrule
\end{tabularx}
\sourcenote{CF-03；Base/Growth为分析拆分，不是公司披露的新分部。}
\end{exhibitbox}

该表的估值含义是：2026 Base Growth毛利率16.8\%隐含显著改善，必须由正式分产品收入与毛利验证。若半年报核心毛利仅维持或下降，模型中最主要的盈利杠杆会失效，即使利润总额仍处预告区间，也可能来自时点性项目或归属变化。

\section{E12｜利润到现金：预付款缓冲与生产占款并存}

2025年末货币资金1,468.35亿元、合同负债1,522.14亿元，看似现金充裕；但存货749.21亿元、合同资产179.78亿元也处高位。合同负债是客户预付款代理，不能等同利润或自由现金；货币资金可能含受限或经营所需资金，也不机械加回股权价值。\srcid{CF-03}

\begin{exhibitbox}[营运资金与现金质量地图]
\small
\begin{tabularx}{\exhibitboxwidth}{L{3.0cm}R{2.5cm}R{2.5cm}L{3.4cm}X}
\toprule
\textbf{指标} & \textbf{2025A} & \textbf{2026Q1} & \textbf{经济含义} & \textbf{风险信号} \\
\midrule
货币资金 & 1,468.35 & 1,445.94 & 生产与交付缓冲 & 不等同可分配净现金 \\
合同负债 & 1,522.14 & 1,523.60 & 客户预付款与订单融资 & 交付后需转收入，不能重复计值 \\
存货 & 749.21 & \na & 在制品、设备与材料投入 & 快于收入增长表明占款加重 \\
合同资产 & 179.78 & \na & 已履约未结算金额 & 增长但收入/回款不跟随是警讯 \\
现金资本开支 & 34.45 & 7.96 & 船位、设备与长期资产投入 & 高景气期扩张可能压低自由现金 \\
经营现金流 & 77.67 & 81.52 & 利润现金化的核心尺度 & 全年OCF/归母低于0.8倍触发降级 \\
\bottomrule
\end{tabularx}
\sourcenote{CF-03、CF-04。金额单位：亿元；Q1缺失字段不以估算补齐。}
\end{exhibitbox}

现金地图的“所以呢”是：客户垫资降低外部借款需求，却不消除生产现金占用。2025短期/长期借款75.76/149.40亿元明显低于货币资金，但这不支持把全部现金作为净现金加回；船舶建造的长周期、合同履约和质保责任要求保留营运资金缓冲。

\section{基础现金情景与质量门槛}

冻结模型预计2026--2028 Base OCF为146.85/205.67/246.81亿元，OCF/归母为0.85/0.95/1.00倍；现金资本开支为50/55/60亿元，对应简化自由现金流96.85/150.67/186.81亿元。该预测不是公司指引，主要用于约束估值倍数是否保留。

\begin{exhibitbox}[Base现金转化路径]
\small
\begin{tabularx}{\exhibitboxwidth}{L{3.4cm}R{2.6cm}R{2.6cm}R{2.6cm}X}
\toprule
\textbf{指标} & \textbf{2026E} & \textbf{2027E} & \textbf{2028E} & \textbf{含义} \\
\midrule
归母净利润 & 172.76 & 216.49 & 246.81 & 合并基础情景 \\
OCF/归母 & 0.85倍 & 0.95倍 & 1.00倍 & 随交付推进逐步改善 \\
经营现金流 & 146.85 & 205.67 & 246.81 & 估值倍数质量门槛 \\
现金资本开支 & 50.00 & 55.00 & 60.00 & 分析假设 \\
简化自由现金流 & 96.85 & 150.67 & 186.81 & OCF减现金资本开支 \\
\bottomrule
\end{tabularx}
\sourcenote{冻结增长模型；金额单位：亿元。自由现金流为简化研究口径。}
\end{exhibitbox}

如果2026全年OCF/归母低于0.8倍且没有可验证的次期回收路径，基础倍数必须下调；若利润强、现金更强且存货与合同资产周转改善，则可提高盈利质量置信度。现金不是估值附录，而是区分“订单兑现”与“会计利润前置”的必要证据。

\begin{riskbox}[半年报现金检查]
正式H1若只兑现利润上端，却出现OCF转负、存货和合同资产快于收入增长，应维持或下调评级，不能用“船舶行业季节性”一语带过。只有披露出清晰的交付回收路径，才能保留2027年正常化倍数。
\end{riskbox}

本章结论是：2025--2026Q1利润弹性已获得披露支持，现金质量仍需完整半年度与全年验证。当前价格接近基础价值时，利润与现金不同步的风险比利润小幅高于预告中点更重要。
